\documentclass[10pt,a4paper]{amsart}
\usepackage{amsmath,amssymb,amsthm,graphicx}
\usepackage{verbatim}
\newenvironment{reference}{\comment}{\endcomment}
%\newenvironment{reference}{}{}
\newenvironment{slogan}{\comment}{\endcomment}
\newenvironment{history}{\comment}{\endcomment}

% For commutative diagrams we use Xy-pic
\usepackage[all]{xy}

% We use 2cell for 2-commutative diagrams.
\xyoption{2cell}
% Theorem environments.
%
\theoremstyle{plain}
\newtheorem{theorem}[subsection]{Theorem}
\newtheorem{proposition}[subsection]{Proposition}
\newtheorem{lemma}[subsection]{Lemma}

\theoremstyle{definition}
\newtheorem{definition}[subsection]{Definition}
\newtheorem{example}[subsection]{Example}
\newtheorem{exercise}[subsection]{Exercise}
\newtheorem{situation}[subsection]{Situation}

\theoremstyle{remark}
\newtheorem{remark}[subsection]{Remark}
\newtheorem{remarks}[subsection]{Remarks}

\numberwithin{equation}{subsection}

% Macros
%
\def\lim{\mathop{\mathrm{lim}}\nolimits}
\def\colim{\mathop{\mathrm{colim}}\nolimits}
\def\Spec{\mathop{\mathrm{Spec}}}
\def\Hom{\mathop{\mathrm{Hom}}\nolimits}
\def\Ext{\mathop{\mathrm{Ext}}\nolimits}
\def\SheafHom{\mathop{\mathcal{H}\!\mathit{om}}\nolimits}
\def\SheafExt{\mathop{\mathcal{E}\!\mathit{xt}}\nolimits}
\def\Sch{\mathit{Sch}}
\def\Mor{\mathop{\mathrm{Mor}}\nolimits}
\def\Ob{\mathop{\mathrm{Ob}}\nolimits}
\def\Sh{\mathop{\mathit{Sh}}\nolimits}
\def\NL{\mathop{N\!L}\nolimits}
\def\CH{\mathop{\mathrm{CH}}\nolimits}
\def\proetale{{pro\text{-}\acute{e}tale}}
\def\etale{{\acute{e}tale}}
\def\QCoh{\mathit{QCoh}}
\def\Ker{\mathop{\mathrm{Ker}}}
\def\Im{\mathop{\mathrm{Im}}}
\def\Coker{\mathop{\mathrm{Coker}}}
\def\Coim{\mathop{\mathrm{Coim}}}

% Boxtimes
%
\DeclareMathSymbol{\boxtimes}{\mathbin}{AMSa}{"02}

%
% Macros for moduli stacks/spaces
%
\def\QCohstack{\mathcal{QC}\!\mathit{oh}}
\def\Cohstack{\mathcal{C}\!\mathit{oh}}
\def\Spacesstack{\mathcal{S}\!\mathit{paces}}
\def\Quotfunctor{\mathrm{Quot}}
\def\Hilbfunctor{\mathrm{Hilb}}
\def\Curvesstack{\mathcal{C}\!\mathit{urves}}
\def\Polarizedstack{\mathcal{P}\!\mathit{olarized}}
\def\Complexesstack{\mathcal{C}\!\mathit{omplexes}}
% \Pic is the operator that assigns to X its picard group, usage \Pic(X)
% \Picardstack_{X/B} denotes the Picard stack of X over B
% \Picardfunctor_{X/B} denotes the Picard functor of X over B
\def\Pic{\mathop{\mathrm{Pic}}\nolimits}
\def\Picardstack{\mathcal{P}\!\mathit{ic}}
\def\Picardfunctor{\mathrm{Pic}}
\def\Deformationcategory{\mathcal{D}\!\mathit{ef}}
\usepackage{cleveref}
\crefname{originallemma}{original lemma}{original lemmas}
\Crefname{originallemma}{Original lemma}{Original lemmas}
\makeatletter
\newlabel{etale-lemma-characterize-etale-completions}{{039M}{source}{}{sourcepointer-1}{}}
\newlabel{etale-lemma-characterize-etale-completions@cref}{{[originallemma][0][]039M}{[1][0][]source}{}{}{}}
\makeatother
\crefname{originalreference}{original reference}{original references}
\Crefname{originalreference}{Original reference}{Original references}
\makeatletter
\newlabel{lemma-node-field-extension}{{0C56}{source}{}{sourcepointer-2}{}}
\newlabel{lemma-node-field-extension@cref}{{[originalreference][0][]0C56}{[1][0][]source}{}{}{}}
\makeatother
\crefname{originallemma}{original lemma}{original lemmas}
\Crefname{originallemma}{Original lemma}{Original lemmas}
\makeatletter
\newlabel{more-algebra-lemma-flat-unramified}{{0AGX}{source}{}{sourcepointer-3}{}}
\newlabel{more-algebra-lemma-flat-unramified@cref}{{[originallemma][0][]0AGX}{[1][0][]source}{}{}{}}
\makeatother
\title{Nodes are stable under etale morphisms}
\author{The Stacks Project Authors}
\date{stacks-project:curves.tex:local-snapshot-20261012; GNU FDL 1.2 or later}
\begin{document}
\maketitle
\noindent\textit{Editorial note.} All delivered mathematical text is the author's own, copied from the pinned Stacks Project chapter file sources/stacks-curves/pinned-original-source.tex; no statement or proof step was rewritten, reordered or removed.\par\medskip
\noindent\textit{Editorial note.} The counted proof is the author's complete argument: the reduction step invokes \textbackslash{}label\{lemma-node-field-extension\} (source line 4669), a published lemma of the same chapter whose own proof delegates its central ring-isomorphism step to the reader (source lines 4646-4652); that lemma was therefore not carded separately, but it is used here as published input per the citation policy; the completion isomorphism and the final conclusion are complete in the present proof.\par\medskip
\noindent\textit{Editorial note.} The label lemma-node-etale-local names the counted statement, and the proof environment immediately following it in the source is delivered verbatim as its proof.\par\medskip
\noindent\textit{Editorial note.} The counted statement is carried by the author's own lemma environment in the chapter source.\par\medskip
\noindent\textit{Allowed standard prerequisites (editorial).}\begin{itemize}
\item Nodes under field extension. Same chapter, Lemma \textbackslash{}label\{lemma-node-field-extension\}. Proof lines 4669-4671: one may base change to the algebraic closure of k, where the residue fields of x and y are both k.
\item Completions of etale local rings. Etale Morphisms, Lemma \textbackslash{}label\{etale-lemma-characterize-etale-completions\}; More on Algebra, Lemma \textbackslash{}label\{more-algebra-lemma-flat-unramified\}. Proof lines 4673-4678: the map on completed local rings O\_\{X,x\}\textasciicircum{} -> O\_\{Y,y\}\textasciicircum{} is an isomorphism.
\item The definition of a node via the completed local ring. Same chapter, Lemma \textbackslash{}label\{lemma-nodal\}, Lemma \textbackslash{}label\{lemma-nodal-algebraic\}. Proof lines 4678-4679: being a node is expressed through the completed local ring (the normal form kappa[[s, t]]/(st)), so an isomorphism of completions preserves it.
\end{itemize}\medskip
\ifdefined\hypertarget\hypertarget{sourcepointer-1}{}\fi
\noindent\textit{Source reference.} Original lemma 039M: sources/stacks-chow/tags.txt, source lines 4277--4277. Prerequisite pointer invoked at proof line(s) 4678; it names a standard fact of a sibling chapter that the counted argument uses, and every step of the counted proof built on it is retained verbatim in the delivered body.\par\medskip
\ifdefined\hypertarget\hypertarget{sourcepointer-2}{}\fi
\noindent\textit{Source reference.} Original reference 0C56: sources/stacks-curves/pinned-original-source.tex, source lines 4605--4605. Prerequisite pointer invoked at proof line(s) 4673; it names a standard fact of the same chapter that the counted argument uses, and every step of the counted proof built on it is retained verbatim in the delivered body.\par\medskip
\ifdefined\hypertarget\hypertarget{sourcepointer-3}{}\fi
\noindent\textit{Source reference.} Original lemma 0AGX: sources/stacks-chow/tags.txt, source lines 13210--13210. Prerequisite pointer invoked at proof line(s) 4679; it names a standard fact of a sibling chapter that the counted argument uses, and every step of the counted proof built on it is retained verbatim in the delivered body.\par\medskip
\begin{lemma}
\label{lemma-node-etale-local}
Let $k$ be a field. Let $X$ be a locally algebraic
$k$-scheme of dimension $1$. Let $Y \to X$ be an \'etale morphism.
Let $y \in Y$ be a point with image $x \in X$. The following are equivalent
\begin{enumerate}
\item $x$ is a closed point of $X$ and a node, and
\item $y$ is a closed point of $Y$ and a node.
\end{enumerate}
\end{lemma}

\begin{proof}
By Lemma \ref{lemma-node-field-extension}
we may base change to the algebraic closure of $k$.
Then the residue fields of $x$ and $y$ are $k$.
Hence the map $\mathcal{O}_{X, x}^\wedge \to \mathcal{O}_{Y, y}^\wedge$
is an isomorphism (for example by
\'Etale Morphisms, Lemma \ref{etale-lemma-characterize-etale-completions} or
More on Algebra, Lemma \ref{more-algebra-lemma-flat-unramified}).
Thus the lemma is clear.
\end{proof}

\end{document}
